\chapter{Diskussion}\label{diskussion}

Die Evaluation in Kapitel~\ref{evaluation} hat die empirische Grundlage gelegt. Dieses Kapitel interpretiert die Evaluationsergebnisse entlang sechs Dimensionen: Diskussion der Betriebsmodi (Abschnitt~\ref{diskussion-der-betriebsmodi}), Diskussion der Prompt-Architektur (Abschnitt~\ref{diskussion-der-prompt-architektur}), Einordnung in den Stand der Technik (Abschnitt~\ref{einordnung-in-den-stand-der-technik}), kritische Reflexion und Limitationen (Abschnitt~\ref{kritische-reflexion}), Implikationen für den Sportjournalismus (Abschnitt~\ref{implikationen-für-den-sportjournalismus}) sowie ethische und regulatorische Einordnung (Abschnitt~\ref{ethik}).


\section{Diskussion der Betriebsmodi}\label{diskussion-der-betriebsmodi}

Die drei Betriebsmodi (\texttt{auto}, \texttt{coop}, \texttt{manual}) wurden in Abschnitt~\ref{status-lifecycle-und-betriebsmodi} konzipiert und in Abschnitt~\ref{evaluation-der-betriebsmodi} evaluiert. Ihre komparative Bewertung ergibt eine klare Hierarchie:

Der \textbf{\texttt{auto}-Modus} maximiert das Latenzpotenzial, trägt jedoch das in Abschnitt~\ref{qualitative-analyse-der-generierten-texte} dokumentierte unkontrollierte Halluzinationsrisiko, das im journalistischen Kontext die Glaubwürdigkeit gefährden kann (Beils 2023, S.~57). Eine Halluzinationsrate von 5\,\% ist dabei im Kontext aktueller LLM-Forschung einzuordnen: Maynez et al.\ (2020, S.~4582) dokumentieren in NLG-Systemen generell substanzielle Raten an Faktinkonsistenzen, insbesondere bei Texten mit Zahlenangaben und Eigennamen. Die Pre-Match-Schutzregel (vgl.~Abschnitt~\ref{prompt-engineering-strategie}) reduziert die kritischste Fehlerklasse (Erfindung von Spielszenen), sodass die verbleibenden 5\,\% ein für LLMs typisches Residualrisiko darstellen, das im \texttt{coop}-Modus durch redaktionelle Sichtung aufgefangen wird. Er ist daher nur für unkritische Event-Typen wie Phasenwechsel vertretbar, bei denen eine fehlerhafte Formulierung keine Falschinformation über Spielergebnisse oder Akteure erzeugt. Darüber hinaus setzt ein verantwortbarer \texttt{auto}-Einsatz organisatorische Rahmenbedingungen voraus: eine nachgelagerte Qualitätsdurchsicht generierter Einträge nach dem Spiel sowie die technische Möglichkeit, fehlerhafte Einträge nachträglich zu korrigieren oder zu löschen. Ohne diese Absicherung entspricht der \texttt{auto}-Modus einer unkontrollierten KI-Publikation ohne redaktionelle Verantwortungslinie.

Der \textbf{\texttt{manual}-Modus} entspricht dem Status quo der manuellen Liveticker-Erstellung und dient im System als Vergleichsbasis sowie als strategische Einstiegsstufe. Er ist für Redaktionen sinnvoll, die das System zunächst parallel zum bestehenden Workflow evaluieren möchten, ohne KI-Inhalte zu publizieren. Darüber hinaus bleibt er als Rückfalloption in Situationen relevant, in denen die LLM-Verfügbarkeit eingeschränkt ist oder das Halluzinationsrisiko für ein spezifisches Spielereignis (z.\,B. strittiger Elfmeter, VAR-Entscheidung) eine vollständig manuelle Einordnung erfordert.

\textbf{Der \texttt{coop}-Modus ist der intendierte Produktivbetrieb}: Die KI liefert Entwürfe, die redaktionelle Freigabe via TAB/ESC-Steuerung sichert Qualität ohne merkliche Latenzzunahme. Methodisch ist dabei zu beachten, dass die gemessene Systemlatenz (Median 3\,338~ms, vgl. Abschnitt~\ref{llm-latenz}) die reine LLM-Generierungszeit erfasst. Im \texttt{coop}-Modus addiert sich die redaktionelle Prüfzeit des Redakteurs, sodass die effektive Time-to-Publish im Produktivbetrieb höher liegt als die gemessene Systemlatenz allein suggeriert. Das Experteninterview (vgl. Abschnitt~\ref{experteninterview-auswertung}) bestätigt diese Einschätzung. Darüber hinaus liefert das Interview ein strukturelles Argument, das über die gemessene Halluzinationsrate hinausgeht. Bei unvorhergesehenen Ereignissen im Stadion, etwa einem Notarzteinsatz während einer 4:0-Führung, würde das System ohne menschliche Prüfung weiterhin euphorisch schreiben, da es nur den Spielstand kennt, nicht die Atmosphäre im Stadion. Diese situative Kontextblindheit (vgl. Abschnitt~\ref{experteninterview-auswertung}) begründet strukturell, warum der \texttt{coop}-Modus auch bei technisch einwandfreier Generierung dem \texttt{auto}-Modus vorzuziehen ist. Die Laufzeit-Umschaltbarkeit ermöglicht eine schrittweise Einführung von \texttt{manual} über \texttt{coop} bis \texttt{auto}, eine Adoptionsstrategie, die dem in Abschnitt~\ref{zeitdruck-und-fehleranfaelligkeit} beschriebenen Akzeptanzbedarf Rechnung trägt. Die modus-spezifische Gesamtbewertung mit Empfehlung findet sich in Abschnitt~\ref{beantwortung-der-forschungsfrage}.


\section{Diskussion der
Prompt-Architektur}\label{diskussion-der-prompt-architektur}

Die Prompt-Architektur verdient eine eigenständige Diskussion, weil ihre Designentscheidungen sowohl die Betriebsmodiwahl (Abschnitt~\ref{diskussion-der-betriebsmodi}) als auch die Evaluationsbefunde aus Abschnitt~\ref{einfluss-der-stilprofile} kausal erklären. Sie gliedert sich in die dynamischen Stilreferenzen (Abschnitt~\ref{stilreferenzen}) und den Temperaturparameter und Context-Builder (Abschnitt~\ref{temperatur-context}).

\subsection{Stilreferenzen}\label{stilreferenzen}

Die Stilkonditionierung beruht auf zwei Schichten. Die primäre Schicht ist \texttt{STYLE\_DESC} (definiert in \texttt{constants.py}): eine statische, profilspezifische Instruktion (50--100~Wörter je Profil), die dem LLM Register, Tonalität, verbotene Floskeln und situationsspezifische Regeln vorgibt. Die sekundäre Schicht sind die dynamisch aus der \texttt{style\_references}-Datenbanktabelle geladenen Beispieltexte (bis zu drei Einträge pro Event-Typ und Instanz), die konkrete Formulierungsbeispiele aus der redaktionellen Praxis liefern. Diese Trennung ermöglicht es, die Stilinstruktion codebasiert zu pflegen und Beispieltexte ohne Codeänderung über die Datenbank zu ergänzen.

Ein strukturell relevanter Unterschied zwischen den Profilen: Die \texttt{euphorisch}-Instruktion enthält explizit die Aufforderung, sich am Rhythmus und der Wortwahl der Stilreferenzen zu orientieren. Sie ist damit auf das Zusammenspiel beider Schichten ausgelegt. Die \texttt{neutral}- und \texttt{kritisch}-Profile verweisen in ihrer Instruktion nicht auf die Beispieltexte und nutzen primär die statische Vorgabe. Diese Asymmetrie hat Konsequenzen für die Qualitätsstabilität, die in Abschnitt~\ref{inhaltliche-reflexion} eingeordnet werden.

Die Referenzen werden nicht nach dem angeforderten Stilprofil gefiltert: Alle drei Profile erhalten dieselben Beispiele aus der Eintracht-Redaktionspraxis, die tendenziell euphorischen Charakter tragen. Dieser Designentscheid erhält den verfügbaren Stichprobenumfang je Bucket, erklärt aber zugleich das in Abschnitt~\ref{einfluss-der-stilprofile} dokumentierte Durchschlagen euphorischer Muster auf das neutrale Profil. Ein kontrollierter A/B-Test (mit vs.~ohne Referenzen, stilgefiltert vs.~stilübergreifend) war nicht durchführbar (vgl. Abschnitt~\ref{einfluss-der-stilprofile}). Qualitativ zeigt die Evaluation konsistente Formatierungsmuster in der \texttt{ef\_whitelabel}-Instanz, insbesondere bei TOOOOR-Konvention und Minutenformat. Der Effekt ist sichtbar, aber statistisch nicht isoliert messbar.

\subsection{Temperaturparameter und Context-Builder}\label{temperatur-context}

Die gewählte Temperatur von 0{,}5 für die Textgenerierung und 0{,}1 für Übersetzungen ist ein Kompromiss zwischen sprachlicher Varianz und Vorhersagbarkeit. Eine niedrigere Temperatur dürfte die Halluzinationsrate weiter senken, die Texte aber formelhafter machen, ein Zielkonflikt, der im Genre Liveticker (das von Überraschung und sprachlicher Kreativität lebt) besonders relevant ist. Die sieben spezialisierten Context-Builder (vgl. Abschnitt~\ref{prompt-engineering-strategie}) strukturieren die Fakten event-typspezifisch vor der Übergabe an das LLM. Die Spezialisierung ist eine bewusste Entscheidung gegen einen generischen Builder: Ein einheitlicher Kontext würde alle verfügbaren Metadaten übergeben (Lineup, H2H-Statistiken, Spielstandverlauf, Spielerprofile) unabhängig vom Event-Typ. Für ein Torereignis sind H2H-Daten irrelevant; für einen Pre-Match-Eintrag sind Minutenangaben bedeutungslos. Irrelevanter Kontext zwingt das Modell zur Selektion und erhöht die Halluzinationswahrscheinlichkeit. Die Builder-Spezialisierung minimiert den Prompt-Umfang auf das event-typisch Notwendige und ist damit eine Anti-Halluzinations-Maßnahme auf Architekturebene. Die Wirksamkeit zeigt sich insbesondere bei ereignistypischen Einträgen (Tore, Karten), wo die Faktengrundlage eng begrenzt ist. Bei Pre-Match-Einträgen verhinderte die Schutzregel zuverlässig Spielszenen-Halluzinationen, nicht jedoch alle Formen der Kontextüberschreitung (vgl. Abschnitt~\ref{qualitative-analyse-der-generierten-texte}).


\section{Einordnung in den Stand der
Technik}\label{einordnung-in-den-stand-der-technik}

Das entwickelte System lässt sich in den Stand der Technik in zwei Richtungen einordnen: die Abgrenzung zu bestehenden Systemen (Abschnitt~\ref{abgrenzung-zu-bestehenden-systemen}) und den Beitrag zur Forschung (Abschnitt~\ref{beitrag-zur-forschung}).

\subsection{Abgrenzung zu bestehenden
Systemen}\label{abgrenzung-zu-bestehenden-systemen}

Die Evaluation ermöglicht eine empirisch gestützte Positionierung des Systems gegenüber bestehenden Ansätzen der automatisierten Sportberichterstattung (Abschnitt~\ref{automatisierte-sportberichterstattung}):

\begin{itemize}\tightlist
\item \textbf{Gegenüber templatebasierten Systemen} (z.\,B. Retresco, vgl. Beils 2023, S.~207) wurde kein direkter empirischer Vergleich durchgeführt. Strukturell zeigt die Evaluation jedoch, dass drei distinkte Stilprofile mit deutlich unterscheidbaren Texten realisierbar sind (neutral Ø~4{,}82, kritisch Ø~4{,}32, euphorisch Ø~3{,}27 im LLM-as-Judge-Verfahren) -- eine adaptive Stilsteuerung, die ein regelbasiertes Template-System ohne manuelle Erweiterung jedes Stilprofils nicht leisten kann. Die Fehlerklasse Phrasenwiederholung (13\,\% der Einträge, vgl. Abschnitt~\ref{qualitative-analyse-der-generierten-texte}) zeigt jedoch, dass auch LLMs ohne gezielte Diversitätsmechanismen in formelhaften Mustern verharren können (eine Schwäche, die templatebasierte Systeme und LLM-basierte Ansätze teilen).
\item \textbf{Gegenüber vollautonomen KI-Systemen} validiert die Evaluation die Entscheidung für den hybriden \texttt{coop}-Modus: Eine Halluzinationsrate von 5\,\% und eine Stil-Inkonsistenzrate von 20\,\% im \texttt{auto}-Modus (vgl. Abschnitt~\ref{qualitative-analyse-der-generierten-texte}) belegen, dass vollautonome Publikation im journalistischen Kontext ein nicht akzeptables Restrisiko trägt (vgl. Bluhm \& Schäfer 2023, S.~33). Das Experteninterview (vgl. Abschnitt~\ref{experteninterview-auswertung}) bestätigt diese Einschätzung: Der Interviewpartner bewertete den \texttt{coop}-Modus als realistischen Kompromiss, der Zeitdruck reduziert, ohne die redaktionelle Kontrolle aufzugeben.
\item \textbf{Gegenüber reinen Übersetzungslösungen} für Mehrsprachigkeit bietet das System nativ multilinguale Generierung durch Prompt-Parametrisierung. Die Übersetzungsqualität wurde jedoch nicht qualitativ evaluiert (vgl. Abschnitt~\ref{technische-einordnung}): Der empirische Abstand zu spezialisierten Übersetzungslösungen (z.\,B. DeepL, NLLB) bleibt damit offen und ist als priorisierter Evaluationsbedarf in Abschnitt~\ref{mittelfristige-erweiterungen} ausgewiesen.
\end{itemize}

\subsection{Beitrag zur Forschung}\label{beitrag-zur-forschung}

Im Sinne der DSR-Artefakttaxonomie nach Hevner et al.\ (2004, S.~77) stellt das System eine \textbf{Instantiierung} dar: ein lauffähiges, prototypisches System, das ein theoretisches Konzept (hybride KI-gestützte Texterstellung mit Human-in-the-Loop) in einer realen Organisationsumgebung realisiert und damit belegt, dass dieses Konzept nicht nur theoretisch formulierbar, sondern praktisch implementierbar und evaluierbar ist. Diese Artefaktkategorie grenzt sich von Konstrukten (begriffsbildenden Taxonomien), Modellen (Referenzarchitekturen ohne Implementierung) und Methoden (Vorgehensmodellen) ab.

Die Arbeit leistet drei Beiträge zum aktuellen Forschungsstand.

Als \textbf{architektonischen Beitrag} demonstriert sie eine produktionsnahe Referenzarchitektur für KI-gestützte Echtzeit-Textgenerierung: Die modulare Provider-Abstraktion mit Prioritätskette, das Few-Shot-Prompting aus einer kuratierten Stilreferenz-Datenbank und das dreistufige Betriebsmodus-Konzept bilden zusammen ein generisch übertragbares Muster für domänenspezifische LLM-Redaktionssysteme. Vergleichbare Architekturen in der Literatur (z.\,B. Graefe 2016; Clerwall 2014) beschreiben überwiegend templatebasierte Systeme zur automatisierten Textgenerierung (Natural Language Generation, NLG) ohne adaptive Stilsteuerung. Ein LLM-basierter Ansatz mit konfigurierbarer Few-Shot-Stilkonditionierung und Human-in-the-Loop-Kontrolle für Liveticker ist nach aktuellem Recherchstand (vgl. Abschnitt~\ref{automatisierte-sportberichterstattung}) eine eigenständige Erweiterung.

Als \textbf{methodischen Beitrag} stellt die Evaluationsinfrastruktur mit Bulk-Generierungsendpunkt, Provider-Override und LLM-as-Judge-Pipeline systematische und reproduzierbare Modellvergleiche bereit, ohne manuelle Einzelbewertung für jede Konfigurationsänderung. Dies löst ein verbreitetes Problem der NLG-Evaluation: die hohe Aufwand-Nutzen-Ratio manueller Bewertungen bei sich ändernden Modellen und Prompts (vgl. Gatt \& Krahmer 2018).

Als \textbf{praktischen Beitrag} greift die White-Label-Architektur mit instanzspezifischer Stilsteuerung den in Abschnitt~\ref{white-label-bedarf-im-profisport} beschriebenen Bedarf von Vereinen als eigenständige Medienproduzenten auf. Die Trennung von generischer Instanz (\texttt{generic}) und vereinsspezifischer Ausprägung (\texttt{ef\_whitelabel}) auf Datenbank- und Prompt-Ebene zeigt, dass ein einziges System mehrere redaktionelle Stimmen bedienen kann, ein Skalierungsmuster, das über den Fußball hinaus anwendbar ist.

Im Sinne der DSR-Evaluationskriterien nach Hevner et al.\ (2004, S.~82--84) erfüllt das Artefakt alle drei Anforderungen: Relevanz durch die direkte Kooperation mit einem Praxispartner im Profifußball (Stackwork GmbH\,/\,Eintracht Frankfurt), methodische Strenge durch eine reproduzierbare, dreistufige Evaluationsinfrastruktur sowie gestalterische Innovation durch die Kombination von Provider-Fallback-Architektur, Few-Shot-Stilkonditionierung und Human-in-the-Loop-Betriebsmoduskonzept.


\section{Kritische Reflexion und Limitationen}\label{kritische-reflexion}

Dieses Kapitel benennt die methodischen, technischen und inhaltlichen Grenzen des Systems und der Evaluation. Es gliedert sich in die methodische Einordnung (Abschnitt~\ref{methodische-einordnung}), die technische Einordnung (Abschnitt~\ref{technische-einordnung}) und die inhaltliche Reflexion (Abschnitt~\ref{inhaltliche-reflexion}).

\subsection{Methodische Einordnung}\label{methodische-einordnung}

Die zentrale methodische Einschränkung betrifft die Selbstbewertung der
Textqualität: Eine unabhängige Bewertung durch professionelle Sportredakteure fehlt. In frühen Iterationszyklen des Design Science Research (DSR), deren Ziel explorative Artefaktvalidierung ist, stellt die Entwickler-Selbstevaluation eine pragmatisch begründete Vorgehensweise dar, die mit zunehmender Artefaktreife durch externe Domänenevaluation zu ergänzen ist (vgl. Hevner et al.~2004, S.~87). Sie begrenzt jedoch die externe Validität der qualitativen Ergebnisse.

Dieser Einschränkung begegnet die Arbeit auf zwei Ebenen: methodisch durch das LLM-as-Judge-Verfahren (N~=~80, Abschnitt~\ref{erweiterte-evaluation}), bei dem ein unabhängiges Modell (\texttt{anthropic/claude-sonnet-4-5}) als Bewerter fungiert und die Selbstreferenzialität systematisch aufbricht. Infrastrukturell wird die implementierte Cohen's-Kappa-Metrik als Vorbereitung für eine Folgestudie mit externen Ratern genutzt. Da das LLM-as-Judge-Verfahren seinerseits modellspezifische Bewertungstendenzen einbringen kann (vgl. Zheng et al.~2023), ersetzt es die externe Nutzerstudie nicht vollständig, erhöht jedoch die Belastbarkeit der Befunde deutlich über eine reine Selbstbewertung hinaus. Konkret ist für das verwendete Bewertungsmodell (\texttt{claude-sonnet-4-5}) eine bekannte Präferenz für sachlich-neutralen Stil dokumentiert. Dies könnte die im Vergleich auffällig niedrigen Tonalitätswerte des euphorischen Stilprofils (LLM-as-Judge~2{,}93/5 vs. Hauptevaluation~3{,}8/5, vgl. Abschnitt~\ref{einfluss-der-stilprofile}) zumindest teilweise erklären. Ob die Diskrepanz primär tatsächliche Qualitätsmängel des euphorischen Profils oder systematischen Bewerterbias widerspiegelt, lässt sich mit den vorliegenden Daten nicht abschließend trennen und bleibt als Validierungsfrage für eine Folgestudie mit menschlichen Ratern offen.

Die eingeschränkte Stichprobengröße (15 Spiele, 40 Events) relativiert die Generalisierbarkeit der
quantitativen Ergebnisse. Insbesondere Randsituationen (torloses
Unentschieden, Elfmeterschießen, Spielabbrüche) sind in der
Datenbasis strukturell unterrepräsentiert. Eine Erweiterung der Datenbasis ist über den bestehenden Bulk-Generierungsendpunkt ohne Codeänderung möglich.

Das strukturierte Experteninterview (Abschnitt~\ref{experteninterview-auswertung}) ergänzt die quantitativen Metriken um domänenspezifische Praxiseinschätzungen, die durch Scores allein nicht erfassbar sind. Als Einschränkung ist das institutionelle Näheverhältnis zwischen Entwickler und Interviewpartner zu berücksichtigen, da der Experte Mitarbeiter der Partnerorganisation Eintracht Frankfurt ist, was einen Gefälligkeitseffekt nicht vollständig ausschließt.

Auf automatische NLG-Evaluationsmetriken wie BLEU (Bilingual Evaluation Understudy) und COMET (Crosslingual Optimized Metric for Evaluation of Translation) wurde bewusst verzichtet. Beide Metriken setzen Referenztexte voraus, gegen die generierte Texte verglichen werden. Da die Liveticker-Generierung eine kreative Aufgabe ohne definierte korrekte Referenzausgabe ist, liefern n-Gramm-Überlappung (BLEU) oder trainierte Ähnlichkeitsmodelle (COMET) keine validen Qualitätsaussagen. Dies ist ein bekanntes Grundproblem der automatischen NLG-Evaluation bei offenen Generierungsaufgaben (vgl. Gatt \& Krahmer 2018, S.~89). Das LLM-as-Judge-Verfahren stellt die methodisch adäquate Alternative dar.

Die externe Validität der Befunde ist an mehrere Kontextgrenzen gebunden, die explizit zu benennen sind: Die Evaluation basiert auf einem Verein (Eintracht Frankfurt), einem LLM-Provider (Gemini~2.0~Flash), einer Sprache (Deutsch) und einer Sportart (Fußball, Bundesliga) über eine Saison. Qualitätswerte, Latenzprofile und Stilprofil-Wirksamkeit können bei anderen Vereins-Korpora, anderen Modellen oder anderen Sportarten (Handball, Basketball, Tennis) systematisch abweichen. Besonders relevant ist dabei, dass die Few-Shot-Referenzen der \texttt{style\_references}-Tabelle ausschließlich aus der redaktionellen Praxis von Eintracht Frankfurt kuratiert sind und nicht auf andere Vereine übertragen werden können. Ein Einsatz für einen anderen Verein erforderte eine vollständige Neukuratierung der Stilbasis. Die Generalisierung der Architekturprinzipien (Provider-Fallback, dreistufige Betriebsmodi, White-Label-Instanzierung) ist davon unabhängig. Die Generalisierung konkreter Qualitätswerte erfordert Folgestudien in anderen Kontexten.

\subsection{Technische Einordnung}\label{technische-einordnung}

Die in Abschnitt~\ref{llm-latenz} dokumentierten Latenzen (Median 3\,338~ms, P95 3\,698~ms) entsprechen einer Kaltgenerierung ohne aktiven Prompt-Cache. Für den Produktivbetrieb ist zu berücksichtigen, dass nach längerer Inaktivität (z.\,B. zwischen zwei Spielen) mit ähnlichen Generierungszeiten zu rechnen ist, da Render-seitige und Gemini-seitige Caches abgelaufen sein können.

Die fehlende Authentifizierung ist im Kontext eines internen Redaktionswerkzeugs
vertretbar, schließt aber einen offenen Mehrbenutzerbetrieb ohne
Netzwerkabsicherung aus.

Die 17 n8n-Workflows stellen den größten nicht automatisiert getesteten
Systemteil dar und verantworten die gesamte Datenversorgung des Systems.

Die n8n-Instanz ist für den Evaluationszeitraum lokal betrieben und über einen \texttt{ngrok}-Tunnel für eingehende Webhooks erreichbar (vgl. Abschnitt~\ref{deployment-konzept}). Diese Lösung ist im Produktivbetrieb nicht vertretbar: Die Tunnel-Stabilität hängt von einer persistenten Internetverbindung des lokalen Rechners ab, und bei Tunnel-Neustarts müssen alle Webhook-URLs aktualisiert werden. Für einen produktiven Betrieb ist ein Cloud-Hosting der n8n-Instanz erforderlich.

Die Provider-Fallback-Kette (Prioritäten 2--4: Gemini direkt, OpenAI, Anthropic) ist architektonisch vollständig implementiert und wurde im Evaluationszeitraum mangels aktivierter API-Keys nicht produktiv eingesetzt. Ein empirischer Qualitätsvergleich der Provider sowie eine Überprüfung der Fallback-Zuverlässigkeit unter Lastbedingungen stehen damit aus, ein relevanter blinder Fleck für ein System, das Multi-Provider-Redundanz als Kern-Architekturmerkmal ausweist (vgl. Abschnitt~\ref{multi-provider-architektur}).

Das System wurde nicht über einen längeren Zeitraum im Produktivbetrieb evaluiert. Aspekte wie kumulierte API-Kosten, Rate-Limit-Häufigkeit im Regelbetrieb und eine mögliche Qualitätsdrift der generierten Texte über eine vollständige Spielsaison hinweg sind empirisch nicht erfasst.

Die Übersetzungsfunktion (\texttt{translate\_text()}, vgl.\ Abschnitt~\ref{llm-service-prompt-aufbau-und-provider-dispatch}) wurde funktional verifiziert, aber nicht qualitativ evaluiert. Eine Beurteilung durch Muttersprachler oder automatische Metriken (z.\,B. BLEU, COMET) wäre Voraussetzung für den Einsatz in nicht-deutschsprachigen Märkten. Die implementierte Batch-Übersetzung erfasst zudem alle publizierten Einträge unabhängig vom \texttt{source}-Feld, sodass auch manuell verfasste Einträge in der Zielsprache übersetzt werden.

\subsection{Inhaltliche Reflexion}\label{inhaltliche-reflexion}

Die Evaluationsergebnisse (Abschnitt~\ref{einfluss-der-stilprofile}) zeigen, dass die
Few-Shot-Referenzen das Format zuverlässig konditionieren:
Minutenangaben, TOOOOR-Konvention und Textkürze werden konsistent
übernommen. Die stilistische Lücke ist damit teilweise geschlossen:
Faktentreue und Verständlichkeit profitieren von der strukturierten
Kontextübergabe, während Tonalität als schwächste Dimension verbleibt.
Die häufigste Fehlerklasse, Stil-Inkonsistenz (vgl.
Fehlerklassen-Tabelle in Abschnitt~\ref{qualitative-analyse-der-generierten-texte}), tritt gerade dort auf, wo
euphorische Few-Shot-Muster auf den neutralen Stil durchschlagen.

Die strukturelle Asymmetrie zwischen den Stilprofilen stellt eine methodische Einschränkung der Evaluation dar: Das \texttt{euphorisch}-Profil ist auf das Zusammenspiel von statischer Instruktion und dynamischen DB-Referenzen ausgelegt, \texttt{neutral} und \texttt{kritisch} primär auf die Instruktion allein (vgl. Abschnitt~\ref{stilreferenzen}). Die schwächeren Tonalitätswerte des euphorischen Profils bei Wechsel- und Karten-Events (LLM-as-Judge Tonalität~2{,}93/5, vgl. Abschnitt~\ref{einfluss-der-stilprofile}) könnten damit auf drei nicht trennbare Faktoren zurückzuführen sein: die STYLE\_DESC-Formulierung, die Referenzdichte für diese Event-Typen in der Datenbank oder die fehlende Stilfilterung der Referenzen. Mit den vorliegenden Daten lässt sich der Beitrag der drei Faktoren nicht isolieren.

Der in Abschnitt~\ref{large-language-models} beschriebene Halluzinationseffekt ist im Kontext von
Livetickern besonders kritisch, da fehlerhafte Fakten (falscher
Torschütze, falsches Ergebnis) unmittelbar die Glaubwürdigkeit
zerstören. Die explizite Schutzregel für Pre-Match-Prompts und die
moderat gewählte Temperatur (0{,}5; vgl. Abschnitt~\ref{llm-service-prompt-aufbau-und-provider-dispatch}) sind Gegenmaßnahmen, deren
Wirksamkeit jedoch nur im \texttt{coop}-Modus durch die redaktionelle
Kontrolle vollständig abgesichert ist. Im \texttt{auto}-Modus verbleibt
ein Restrisiko fehlerhafter Veröffentlichungen.

Zugleich ist zu berücksichtigen, dass der hohe Korrektheitswert (Ø~4{,}6\,/\,5) durch die strukturierte Faktenübergabe im Prompt begünstigt wird: Torschütze, Minute und Spielstand stammen direkt aus der API und werden dem Modell explizit übergeben. Ein Modell, das den Prompt nicht ignoriert, erzeugt damit faktuell korrekte Einträge weitgehend unabhängig von seiner eigentlichen Sprachkompetenz. Der Korrektheitswert misst damit eher die Güte des Prompt-Designs als eine intrinsische LLM-Stärke, was die hohen Werte erklärt, deren Übertragbarkeit auf faktenarme Prompts jedoch einschränkt.


\section{Implikationen für den
Sportjournalismus}\label{implikationen-fuxfcr-den-sportjournalismus}

Das hybride System zielt darauf ab, die Rolle des Liveticker-Redakteurs von
der Textproduktion zur Textredaktion zu verlagern: Statt unter Zeitdruck zu
formulieren, prüft, editiert und kuratiert der Redakteur KI-generierte
Vorschläge. Diese Verschiebung entspricht dem in der Forschungsliteratur diskutierten Wandel hin zu einer
„augmented authorship" (erweiterter Autorschaft; vgl. Bluhm \& Schäfer 2023, S.~35),
bei der menschliche Expertise und maschinelle Effizienz komplementär
zusammenwirken. Der kognitive Engpass der simultanen Beobachtung des Spielgeschehens, Relevanzbewertung und Texterstellung unter extremem Zeitdruck (vgl. Abschnitt~\ref{zeitdruck-und-fehleranfaelligkeit}) könnte dadurch wesentlich reduziert werden. Ob dies in der Praxis zu messbaren Effizienzgewinnen führt, ist empirisch nicht belegt. Das Experteninterview (vgl. Abschnitt~\ref{experteninterview-auswertung}) deutet auf Akzeptanz und erwarteten Nutzen hin, ersetzt jedoch keinen Feldtest mit echten Redakteuren unter Live-Bedingungen.
Die originäre Kompetenz würde sich verlagern: Sprachliche
Formulierungsfähigkeit tritt in den Hintergrund, evaluative Kompetenz
(das schnelle Erkennen von Stilbrüchen, Faktenfehlern und
Tonalitätsabweichungen) wird zur Kernkompetenz. Die publizistische
Verantwortung verbleibt vollständig beim Menschen.

Die White-Label-Architektur (\texttt{ef\_whitelabel} vs.~\texttt{generic})
trägt dem strukturellen Wandel der Vereine zu eigenständigen
Medienproduzenten Rechnung (vgl. Abschnitt~\ref{white-label-bedarf-im-profisport}). Ein einzelnes System kann, durch
Instanzkonfiguration, Stilprofile und Few-Shot-Referenzen,
verschiedene redaktionelle Stimmen bedienen, ohne separate Codebases zu
erfordern. Technische Voraussetzung für eine echte Multi-Vereins-Skalierung ist Mandantenfähigkeit auf Datenbankebene. Die aktuelle Architektur hält beide Instanzen in einem Schema. Konkrete Erweiterungsszenarien sind in
Abschnitt~\ref{ausblick} beschrieben.

Die implementierte Mehrsprachigkeit adressiert darüber hinaus die in Abschnitt~\ref{mehrsprachigkeit-als-herausforderung} motivierte Internationalisierungsanforderung: Vereine mit globalen Fanbases könnten KI-generierte Einträge ohne parallele Redaktionskapazitäten in mehreren Sprachen bereitstellen, sofern die Übersetzungsqualität durch eine externe Evaluation abgesichert wird (vgl. Abschnitt~\ref{mittelfristige-erweiterungen}).


\section{Ethische und regulatorische Einordnung}\label{ethik}

Der Einsatz KI-generierter Texte im Sportjournalismus berührt normative Fragen, die über die technische Systemgestaltung hinausgehen. Sie gliedert sich in die ethische Einordnung (Abschnitt~\ref{ethische-einordnung}) und die regulatorischen Anforderungen (Abschnitt~\ref{regulatorische-anforderungen}).

\subsection{Ethische Einordnung}\label{ethische-einordnung}

\paragraph{Transparenz gegenüber Lesern.}
Die Veröffentlichung KI-generierter Liveticker-Einträge ohne Kennzeichnung ist aus journalistisch-ethischer Sicht kritisch zu reflektieren: Leser können nicht zwischen redaktionell verantworteten und automatisch generierten Texten unterscheiden. Der \texttt{auto}-Modus verstärkt dieses Risiko, da Halluzinationen (Tonalitätsinkonsistenz 20\,\%, Halluzinationsrate 5\,\%) ohne menschliche Prüfung publiziert werden können. Aus journalistischer Sicht entspräche eine transparente Kennzeichnung KI-generierter Inhalte dem Gebot der Quellentransparenz (vgl. Bluhm \& Schäfer 2023, S.~57). Technisch wäre eine solche Kennzeichnung über das bestehende \texttt{source}-Feld im Ticker-Datenmodell (\texttt{ai} vs.~\texttt{manual}) ohne Codeänderung realisierbar.

\paragraph{Publizistische Verantwortung.}
Das Human-in-the-Loop-Design des \texttt{coop}-Modus beantwortet die Verantwortungsfrage konstruktiv: Die redaktionelle Freigabe stellt sicher, dass ein Mensch die publizistische Verantwortung für jeden veröffentlichten Eintrag trägt. Im \texttt{auto}-Modus ist diese Verantwortungslinie unterbrochen, eine Einschränkung, die über die technische Halluzinationsrate hinaus eine normative Systementscheidung darstellt.

\subsection{Regulatorische Anforderungen}\label{regulatorische-anforderungen}

\paragraph{Regulatorischer Rahmen.}
Ab August 2026 verpflichtet Art.~50 Abs.~4 des EU AI Act zur Kennzeichnung KI-generierter Texte, sofern die automatisierte Herkunft nicht offensichtlich erkennbar ist. Liveticker-Einträge fallen unter diesen Anwendungsbereich. Ein produktiver Betrieb des Systems erfordert damit eine explizite Kennzeichnungsstrategie sowie datenschutzkonforme Verarbeitungsvereinbarungen mit den LLM-Providern (insbesondere OpenRouter als Intermediär). Diese regulatorischen Anforderungen lagen außerhalb des Projektumfangs, sind aber für eine Produktiveinführung unerlässlich.

Die Diskussion entlang der sechs Dimensionen ergibt ein konsistentes Gesamtbild: Das System bewältigt die identifizierten Problemdimensionen durch nachweisbare Mechanismen. Verbleibende Schwächen (Stil-Inkonsistenz, fehlende Authentifizierung, eingeschränkte Stichprobengröße) sind klar lokalisiert und behebbar. Kapitel~\ref{fazit-und-ausblick} fasst die Erkenntnisse zusammen, beantwortet die Forschungsfrage und gibt einen Ausblick auf Erweiterungen.
